% Abstract — Physics-anchored tandem neural networks for multi-objective
% inverse design of low-cost polymer diffractive AR waveguides
%
% Source of truth for all numbers: docs/ABSTRACT.md ("Where each number
% comes from" section) — regenerate this file by hand if those results
% change (5-seed v3 protocol run, 2026-07-17). Standalone-compilable;
% drop the \documentclass...\begin{document} wrapper when pasting into a
% venue template (Optics Express / SPIE) that supplies its own.

\documentclass[11pt]{article}
\usepackage[utf8]{inputenc}
\usepackage[T1]{fontenc}
\usepackage{amsmath,amssymb}
\usepackage[margin=1in]{geometry}
\usepackage{times}

\title{Physics-Anchored Tandem Neural Networks for Multi-Objective\\
Inverse Design of Low-Cost Polymer Diffractive AR Waveguides}
\author{Connor Wang}
\date{}

\begin{document}
\maketitle

\begin{abstract}
Diffractive waveguide combiners dominate consumer augmented-reality (AR)
displays, yet their design couples material parameters and grating geometry
to system-level image quality in ways that make manual iteration slow, and
existing neural inverse-design approaches invert a single component-level
metric through a learned forward surrogate that itself carries approximation
error. We present a neural inverse-design framework for poly(methyl
methacrylate) (PMMA) surface-relief-grating waveguides that inverts a joint,
perceptually grounded system specification --- modulation transfer function
(MTF) at 40~cyc/mm, double-coupler transmission, lateral chromatic spread,
and transmission at the field edge --- and anchors every stage to exact
physics. A differentiable forward model enforces the total-internal-reflection
guiding condition, resolves Fresnel and grating responses by polarization,
and replaces scalar grating-coupling theory --- which we show overestimates
first-order coupling efficiency $5$--$15\times$ at the total-internal-reflection
(TIR)-mandated sub-wavelength periods ($430$--$449$~nm) --- with a
polarization-resolved efficiency term calibrated to $90{,}090$ rigorous
coupled-wave analysis (RCWA) solutions over refractive index, period,
depth, and duty cycle (mean interpolation error $5.6\times10^{-4}$,
audited against off-grid solves). A forward surrogate
network reproduces this engine with held-out $R^2 \geq 0.9989$ on every
metric across five independent retrainings and passes a four-test
memorization audit (never-seen error $1.2\times$ training error); a tandem
inverse network trained through the frozen forward model resolves the
one-to-many inverse problem, recovering designs that meet requested
specifications in milliseconds with $0.8$--$7\%$ median per-metric error ---
where a direct design-regression baseline fails at $12$--$47\%$, the
signature of inverse non-uniqueness --- and a neural-adjoint search over the
trained surrogate discovers record designs that an independent
exact-physics gradient search reproduces. Rigorous calibration qualitatively
changes the optimum: the best design's grating depth moves from the
scalar-theory bound (400~nm) to an interior optimum at $\approx$200~nm and
its duty cycle from $0.5$ to $\approx$0.41, with transmission quoted at
rigorous rather than scalar accuracy ($1.3\%$ double-coupler, in line with
reported single-digit system efficiencies for diffractive combiners) and MTF
at 40~cyc/mm improving to $0.78$. The framework provides a template for
physics-honest, system-level inverse design of consumer-grade polymer
waveguides.
\end{abstract}

\bigskip
\noindent\textbf{Summary statement.}
Prior tandem-network grating design inverts component-level spectra through
learned RCWA surrogates; this work is, to our knowledge, the first to invert
a joint display-level specification (MTF, transmission, chromatic spread,
field-edge transmission simultaneously) while anchoring the training loss to
an exact differentiable physics engine whose coupling term is calibrated to
rigorous electromagnetic solutions, and the first neural inverse-design
treatment of PMMA --- the lowest-cost, lightest waveguide substrate --- under
its full-RGB TIR feasibility window. Every headline design is re-verified
per-polarization with an independent rigorous solver, and the surrogate's
generalization is demonstrated with an explicit memorization audit, both of
which directly answer the reviewer objections that most commonly meet
surrogate-based inverse design.

\end{document}
